\chapter*{Abstract}

\section*{\tituloPortadaEngVal}

This Bachelor's thesis addresses the design and development of a computer-vision-based tool for the segmentation of magnetic resonance imaging (MRI) images, with the objective of obtaining volumetric measurements from anatomical and pathological regions within a preclinical research framework.

A functional prototype has been developed, in collaboration with ICTS BioImagen Complutense, focused on one of the use cases proposed by the centre: the segmentation of ischaemic lesions in mouse brain. The system uses two segmentation models, applied sequentially to each MRI slice: first, to delineate the whole brain and, subsequently, to identify the damaged region. The volumes of the whole brain and the affected area are calculated from the obtained segmentation masks.

Due to the scarcity of annotated data, which is a common challenge in biomedical imaging, the proposed methodology relies on \textit{transfer learning} within a U-Net architecture. In addition, Explainable Artificial Intelligence (XAI) techniques, such as Grad-CAM and \textit{Integrated Gradients}, have been incorporated to facilitate the qualitative analysis of the predictions and increase the transparency of the prototype.

\section*{Keywords}

\noindent Magnetic Resonance Imaging (MRI), Cerebral ischaemia, Medical segmentation, Region of Interest (ROI), Artificial Intelligence (AI), Convolutional Neural Networks (CNN), U-Net, Explainable Artificial Intelligence (XAI), Computer vision.